% Technical appendix: roadmap paragraph (precedes Appendices A--F).
\section*{\LARGE Technical appendix}
\phantomsection\addcontentsline{toc}{section}{Technical appendix}
\markboth{Technical appendix}{}

The appendices fall into three groups. Each can be read on its own, and every entry, proof and table points back to the
place in the main text that uses it.

\paragraph{Preliminaries}
Appendix~\ref{app:background} collects the mathematical background the exposés assume: categories, slices and base
change; monoidal categories, string diagrams and lenses; manifolds, germs and cones; and groups, symmetry and
conservation laws. Its 93 entries keep the numbers A.1--D.34 they carry on the companion site, and each states a notion
with its hypotheses, shows it on a fine-tuning example and lists the results that rely on it. Appendix~\ref{app:auxiliary}
quotes the external theorems that the proofs use, each in the exact form in which it is applied and with its source.

\paragraph{Theory}
Appendix~\ref{app:proofs} gives complete proofs of the numbered results, in the order of the main text and with one
subsection per result. Each subsection restates its result, in the form recorded after two rounds of review, before the
proof.

\paragraph{Data and reproducibility}
Appendix~\ref{app:catalogue} describes the catalogue of \nMethods{} methods from \nPapers{} papers. It lists every method
with its seven coordinates and its reference, and every one of the \nArrows{} recorded simulation arrows with its map $h$.
Appendix~\ref{app:repro} documents the numerical checks marked [Num] in the main text and explains how to rerun them from
the data files that accompany this paper.

\paragraph{Process}
Appendix~\ref{app:making-of} records how this paper and its companion site were made: who did what, how the literature was searched and checked, how the theory was generalised and refereed, and what the work cost in agent runs and tokens.
